\documentclass[aspectratio=169,10pt]{beamer}
\usetheme[numbering=fraction,progressbar=frametitle]{metropolis}
\usepackage{fontspec}
\setsansfont{Latin Modern Sans}
\usepackage{amsmath,amssymb,mathtools,graphicx,booktabs}
\usepackage{polyglossia}
\setdefaultlanguage{french}
\setbeamertemplate{navigation symbols}{}
\setbeamerfont{frametitle}{size=\large}
\setbeamerfont{normal text}{size=\normalsize}
\setlength{\parskip}{.45em}
\hypersetup{pdfauthor={},pdfsubject={Lecture slides},pdfkeywords={},pdftitle={Fonctions : représentations et fonctions de référence}}
\providecommand{\N}{\mathbb{N}}
\providecommand{\Z}{\mathbb{Z}}
\providecommand{\D}{\mathbb{D}}
\providecommand{\Q}{\mathbb{Q}}
\providecommand{\R}{\mathbb{R}}
\providecommand{\C}{\mathbb{C}}
\providecommand{\vect}[1]{\overrightarrow{#1}}
\providecommand{\vi}{\vec{\imath}}
\providecommand{\vj}{\vec{\jmath}}
\providecommand{\vk}{\vec{k}}
\providecommand{\norm}[1]{\left\lVert #1\right\rVert}
\providecommand{\vc}[2]{\begin{pmatrix}#1\\#2\end{pmatrix}}
\providecommand{\deter}[2]{\det\left(#1,#2\right)}
\providecommand{\abs}[1]{\left\lvert #1\right\rvert}
\providecommand{\dd}{\,\mathrm{d}}
\providecommand{\e}{\mathrm{e}}
\providecommand{\Card}{\operatorname{Card}}
\providecommand{\E}{\operatorname{E}}
\providecommand{\V}{\operatorname{V}}
\providecommand{\degre}{^{\circ}}
\providecommand{\un}[1]{\,\mathrm{#1}}
\providecommand{\sci}[2]{#1\times10^{#2}}
\providecommand{\pgcd}{\operatorname{PGCD}}
\providecommand{\ppcm}{\operatorname{PPCM}}
\title{Fonctions : représentations et fonctions de référence}
\subtitle{Cours et exemples}
\date{}
\author{}
\begin{document}
\begin{frame}[plain]
\vfill
{\usebeamerfont{title}\inserttitle\par}
\medskip
{\usebeamerfont{subtitle}\insertsubtitle\par}
\medskip
{\usebeamercolor[fg]{progress bar}\rule{\textwidth}{.5pt}}
\vfill
\end{frame}
\begin{frame}{Fonction, image et antécédent}
Une fonction à valeurs réelles associe à chaque réel \(x\) de son ensemble de définition \(D_f\) un unique réel \(f(x)\), appelé image de \(x\). Un antécédent de \(y\) est un réel \(x\in D_f\) tel que \(f(x)=y\). On note \(f:x\mapsto f(x)\). Un nombre peut avoir aucun, un ou plusieurs antécédents.
\end{frame}
\begin{frame}{Fonction, image et antécédent}
Soit \(g(x)=x^2-4x+32\) sur \(\R\). Calculer \(g(-5)\). Soit \(h(x)=3x+5\) ; déterminer l'antécédent de \(6\) par \(h\).
\end{frame}
\begin{frame}{Correction : Fonction, image et antécédent}
\(g(-5)=25+20+32=77\) : \(77\) est l'image de \(-5\), et \(-5\) est un antécédent de \(77\). Pour \(h\), \(3x+5=6\iff x=1/3\). L'unique antécédent est \(1/3\).
\end{frame}
\begin{frame}{Fonction, image et antécédent}
Pour la fonction carré, calculer les images de \(-3\) et \(3\), puis les antécédents de \(9\), \(0\) et \(-1\).
\end{frame}
\begin{frame}{Correction : Fonction, image et antécédent}
\(f(-3)=f(3)=9\). Les antécédents de \(9\) sont \(-3,3\) ; \(0\) a un seul antécédent, \(0\) ; \(-1\) n'en a aucun dans \(\R\).
\end{frame}
\begin{frame}[fragile]{Python}
\small
\begin{verbatim}
def g(x):
    return x*x - 4*x + 32
\end{verbatim}
\end{frame}
\begin{frame}{Fonction, image et antécédent}
Distinguer \texttt{g(-5)} de \texttt{x = -5}.
\end{frame}
\begin{frame}{Correction : Fonction, image et antécédent}
Le premier appelle la fonction et renvoie \(77\). Le second affecte \(-5\) à la variable \texttt{x}, sans calculer l'image par \(g\).
\end{frame}
\begin{frame}{Ensemble de définition}
\(D_f\) contient les valeurs pour lesquelles le procédé définissant \(f\) a un sens. Un dénominateur doit être non nul ; un nombre placé sous une racine carrée doit être positif ou nul. Respecter aussi tout domaine imposé par l'énoncé ou le contexte.
\end{frame}
\begin{frame}{Ensemble de définition}
Déterminer les ensembles de définition :
\[f(x)=3x^2-5x+2,\quad g(x)=\dfrac{2x+1}{x-4},\quad h(x)=\dfrac{x}{x^2-9},\quad k(x)=\sqrt{2x-6}.\]
\end{frame}
\begin{frame}{Correction : Ensemble de définition}
\(D_f=\R\). \(D_g=]-\infty;4[\cup]4;+\infty[\) car \(x-4\ne0\). \(D_h=]-\infty;-3[\cup]-3;3[\cup]3;+\infty[\) car \(x^2-9\ne0\). \(D_k=[3;+\infty[\) car \(2x-6\geq0\).
\end{frame}
\begin{frame}{Ensemble de définition}
La racine carrée d'un nombre négatif n'existe pas dans \(\R\) ; ce n'est pas « une racine carrée négative » que l'on exclut.
\end{frame}
\begin{frame}{Ensemble de définition}
Un rectangle de périmètre \(20\) m a pour largeur \(x\) m. Exprimer son aire et préciser le domaine si les deux dimensions doivent être strictement positives.
\end{frame}
\begin{frame}{Correction : Ensemble de définition}
La longueur vaut \(10-x\), donc \(A(x)=x(10-x)\) en m\(^2\). \(x>0\) et \(10-x>0\) donnent \(D_A=]0;10[\). Si les rectangles dégénérés étaient autorisés, on pourrait prendre \([0;10]\).
\end{frame}
\begin{frame}{Tableau de valeurs et courbe représentative}
\(\mathcal C_f\) est l'ensemble des points \((x;f(x))\), avec \(x\in D_f\). Ainsi \(M(a;b)\in\mathcal C_f\iff a\in D_f\) et \(b=f(a)\).
\end{frame}
\begin{frame}{Tableau de valeurs et courbe représentative}
\begin{columns}[T,onlytextwidth]
\begin{column}{.57\textwidth}
Sur \([-1;4]\), \(f(x)=x^2-2x-1\). Compléter le tableau, puis tracer la courbe sur tout l'intervalle.
\end{column}
\begin{column}{.4\textwidth}
\centering
\includegraphics[width=\linewidth,height=4.5cm,keepaspectratio]{figures/seconde-06-1.pdf}
\end{column}
\end{columns}
\end{frame}
\begin{frame}{Correction : Tableau de valeurs et courbe représentative}
Images : \(2\) ; \(0{,}25\) ; \(-1\) ; \(-2\) ; \(-1{,}75\) ; \(-1\) ; \(2\) ; \(7\). Placer les points puis tracer la parabole. Des segments reliant ces seuls points ne donneraient qu'une approximation.
\end{frame}
\begin{frame}{Résolution graphique d'équations et d'inéquations}
Les solutions de \(f(x)=k\) sont les abscisses des intersections de \(\mathcal C_f\) avec \(y=k\). Pour \(f(x)\leq k\), lire les abscisses où la courbe est en dessous de cette droite, intersections comprises ; avec \(<\), les intersections sont exclues.
\end{frame}
\begin{frame}{Résolution graphique d'équations et d'inéquations}
\begin{columns}[T,onlytextwidth]
\begin{column}{.57\textwidth}
La courbe représente \(f(x)=1+\dfrac14(x+3)(x+1)(x-2)\) sur \([-4;3]\). Résoudre graphiquement \(f(x)=1\), \(f(x)\leq1\) et \(f(x)>1\).
\end{column}
\begin{column}{.4\textwidth}
\centering
\includegraphics[width=\linewidth,height=4.5cm,keepaspectratio]{figures/seconde-06-2.pdf}
\end{column}
\end{columns}
\end{frame}
\begin{frame}{Correction : Résolution graphique d'équations et d'inéquations}
\(f(x)=1\) : \(-3,-1,2\). \(f(x)\leq1\) : \([-4;-3]\cup[-1;2]\). \(f(x)>1\) : \(]-3;-1[\cup]2;3]\). Les bornes \(-4,3\) viennent du domaine. La factorisation vérifie exactement les intersections ; l'étude algébrique du signe sera consolidée en séquence 8.
\end{frame}
\begin{frame}{Comparer deux fonctions ; lire un signe}
\(f(x)=g(x)\) se lit aux intersections des courbes. Pour \(f(x)<g(x)\), \(\mathcal C_f\) est strictement en dessous de \(\mathcal C_g\) à la même abscisse. Le signe de \(f(x)\) se lit par rapport à l'axe des abscisses.
\end{frame}
\begin{frame}{Comparer deux fonctions ; lire un signe}
\begin{columns}[T,onlytextwidth]
\begin{column}{.57\textwidth}
Sur \([-3;3]\), \(f(x)=x^2\) et \(g(x)=x+2\). Résoudre \(f=g\), \(f<g\), \(f\geq g\) et déterminer le signe de \(g\).
\end{column}
\begin{column}{.4\textwidth}
\centering
\includegraphics[width=\linewidth,height=4.5cm,keepaspectratio]{figures/seconde-06-3.pdf}
\end{column}
\end{columns}
\end{frame}
\begin{frame}{Correction : Comparer deux fonctions ; lire un signe}
\(f=g\) pour \(x=-1,2\) ; \(f<g\) sur \(]-1;2[\) ; \(f\geq g\) sur \([-3;-1]\cup[2;3]\). \(g\) est strictement négative sur \([-3;-2[\), nulle en \(-2\), strictement positive sur \(]-2;3]\). Signe et variations sont des informations différentes.
\end{frame}
\begin{frame}{Fonctions affines et valeur absolue}
Une fonction affine est \(f(x)=mx+p\) sur \(\R\). Elle est linéaire si \(p=0\), constante si \(m=0\). Pour \(m\ne0\), elle s'annule en \(-p/m\) ; son signe est celui de \(m\) à droite de ce nombre.
La fonction \(x\mapsto|x|\) est définie sur \(\R\), positive ou nulle, et s'annule uniquement en \(0\).
\end{frame}
\begin{frame}{Fonctions affines et valeur absolue}
\begin{columns}[T,onlytextwidth]
\begin{column}{.57\textwidth}
Classer (plusieurs réponses possibles) en affine, linéaire, constante : \(5x+2\) ; \(5x\) ; \(7+2x-7\) ; \(6\) ; \(6(4x-2)\) ; \(6x+5-6x\).
\end{column}
\begin{column}{.4\textwidth}
\centering
\includegraphics[width=\linewidth,height=4.5cm,keepaspectratio]{figures/seconde-06-4.pdf}
\end{column}
\end{columns}
\end{frame}
\begin{frame}{Correction : Fonctions affines et valeur absolue}
Toutes sont affines. \(5x\) et \(2x\) sont linéaires ; \(6\) et \(5\) sont constantes ; \(24x-12\) n'est ni linéaire ni constante.
\end{frame}
\begin{frame}{Fonctions affines et valeur absolue}
Si \(k<0\), \(|x|=k\) et \(|x|\leq k\) n'ont aucune solution, et \(|x|\geq k\) est vraie sur \(\R\). Si \(k=0\), les deux premières donnent \(x=0\), la troisième est vraie sur \(\R\). Pour \(k>0\) :
\[|x|=k\iff x\in\{-k;k\},\quad |x|\leq k\iff x\in[-k;k],\]
\[|x|\geq k\iff x\in]-\infty;-k]\cup[k;+\infty[.\]
\end{frame}
\begin{frame}{Fonctions affines et valeur absolue}
Résoudre \(|x|=5\), \(|x|=-3\), \(|x|\leq3\), \(|x|\geq4\).
\end{frame}
\begin{frame}{Correction : Fonctions affines et valeur absolue}
\(\{-5;5\}\) ; \(\varnothing\) ; \([-3;3]\) ; \(]-\infty;-4]\cup[4;+\infty[\).
\end{frame}
\begin{frame}{Fonctions carré et inverse}
\begin{columns}[T,onlytextwidth]
\begin{column}{.57\textwidth}
\(x\mapsto x^2\) est définie sur \(\R\) ; sa courbe est une parabole. \(x^2\geq0\), avec égalité seulement en \(0\).
\(x\mapsto1/x\) est définie sur \(\R^*\) ; sa courbe est une hyperbole. \(1/x\) a le signe de \(x\) et ne s'annule jamais.
\end{column}
\begin{column}{.4\textwidth}
\centering
\includegraphics[width=\linewidth,height=2.25cm,keepaspectratio]{figures/seconde-06-5.pdf}
\includegraphics[width=\linewidth,height=2.25cm,keepaspectratio]{figures/seconde-06-6.pdf}
\end{column}
\end{columns}
\end{frame}
\begin{frame}{Fonctions carré et inverse}
Si \(k<0\), \(x^2=k\) et \(x^2\leq k\) n'ont pas de solution ; \(x^2\geq k\) est vraie sur \(\R\). Si \(k=0\), les deux premières donnent \(x=0\), la troisième est vraie sur \(\R\). Pour \(k>0\) :
\[x^2=k\iff x\in\{-\sqrt k;\sqrt k\},\quad x^2\leq k\iff x\in[-\sqrt k;\sqrt k],\]
\[x^2\geq k\iff x\in]-\infty;-\sqrt k]\cup[\sqrt k;+\infty[.\]
\end{frame}
\begin{frame}{Fonctions carré et inverse}
Résoudre \(x^2=4\), \(x^2=-1\), \(x^2\geq4\), \(x^2\geq-1\), \(1/x=4\), \(1/x=0\).
\end{frame}
\begin{frame}{Correction : Fonctions carré et inverse}
\(\{-2;2\}\) ; \(\varnothing\) ; \(]-\infty;-2]\cup[2;+\infty[\) ; \(\R\) ; \(\{1/4\}\) ; \(\varnothing\). Plus généralement, \(1/x=k\) a l'unique solution \(1/k\) si \(k\ne0\), aucune si \(k=0\).
\end{frame}
\begin{frame}{Inéquations avec la fonction inverse}
Respecter \(x\ne0\) et lire les deux branches de l'hyperbole. Ne pas multiplier une inéquation par \(x\) sans connaître son signe.
\end{frame}
\begin{frame}{Inéquations avec la fonction inverse}
Résoudre \(1/x\geq2\), \(1/x\leq-1\), \(1/x<1/2\) et \(1/x>-2\).
\end{frame}
\begin{frame}{Correction : Inéquations avec la fonction inverse}
\(]0;1/2]\) ; \([-1;0[\) ; \(]-\infty;0[\cup]2;+\infty[\) ; \(]-\infty;-1/2[\cup]0;+\infty[\). Une inégalité stricte exclut l'intersection avec \(y=k\) ; \(0\) est toujours interdit.
\end{frame}
\begin{frame}{Inéquations avec la fonction inverse}
Un élève écrit : « \(1/x\leq2\iff1\leq2x\), donc \(x\geq1/2\) ». Que manque-t-il ?
\end{frame}
\begin{frame}{Correction : Inéquations avec la fonction inverse}
Il a supposé implicitement \(x>0\). Pour \(x<0\), \(1/x<0<2\) : ces valeurs conviennent aussi. L'ensemble correct est \(]-\infty;0[\cup[1/2;+\infty[\). Le nombre \(-1\) réfute son équivalence.
\end{frame}
\begin{frame}{Mobiliser racine carrée et cube}
\begin{columns}[T,onlytextwidth]
\begin{column}{.57\textwidth}
\(x\mapsto\sqrt x\) est définie sur \([0;+\infty[\) et à valeurs positives ou nulles. Pour \(k\geq0\), \(\sqrt x=k\iff x=k^2\) ; pour \(k<0\), aucune solution.
\(x\mapsto x^3\) est définie sur \(\R\). Chaque réel \(k\) possède un unique antécédent, noté \(\sqrt[3]{k}\).
\end{column}
\begin{column}{.4\textwidth}
\centering
\includegraphics[width=\linewidth,height=2.25cm,keepaspectratio]{figures/seconde-06-7.pdf}
\includegraphics[width=\linewidth,height=2.25cm,keepaspectratio]{figures/seconde-06-8.pdf}
\end{column}
\end{columns}
\end{frame}
\begin{frame}{Mobiliser racine carrée et cube}
Résoudre \(\sqrt x=3\) ; \(\sqrt x=-2\) ; \(\sqrt x\leq2\) ; \(\sqrt x>3\) ; \(x^3=-8\) ; \(x^3\leq27\).
\end{frame}
\begin{frame}{Correction : Mobiliser racine carrée et cube}
\(\{9\}\) ; \(\varnothing\) ; \([0;4]\) ; \(]9;+\infty[\) ; \(\{-2\}\) ; \(]-\infty;3]\). Respecter le domaine pour les racines carrées. La fonction cube conserve l'ordre.
\end{frame}
\begin{frame}{Mobiliser racine carrée et cube}
Comparer \(x^2\) et \(x\) pour \(x=1/2\), puis \(x=2\). Le carré d'un positif est-il toujours plus grand que lui ?
\end{frame}
\begin{frame}{Correction : Mobiliser racine carrée et cube}
\(1/4<1/2\) mais \(4>2\). L'affirmation est fausse ; \(1/2\) est un contre-exemple. Étude générale en séquence 10.
\end{frame}
\end{document}
